% Generated draft: requires mathematical and source review.
\documentclass[11pt]{article}
\usepackage[margin=25mm]{geometry}
\usepackage{fontspec}
\setmainfont{Latin Modern Roman}
\usepackage{amsmath,amssymb,mathtools,unicode-math}
\setmathfont{Latin Modern Math}
\usepackage{enumitem,needspace,xurl,hyperref}
\setlength{\parindent}{0pt}
\setlength{\parskip}{5pt}
\setlength{\emergencystretch}{3em}
\newcommand{\N}{\mathbb{N}}
\newcommand{\Z}{\mathbb{Z}}
\newcommand{\Q}{\mathbb{Q}}
\newcommand{\R}{\mathbb{R}}
\newcommand{\C}{\mathbb{C}}
\newcommand{\F}{\mathbb{F}}
\newcommand{\A}{\mathbb{A}}
\newcommand{\PP}{\mathbb{P}}
\newcommand{\E}{\mathbb{E}}
\newcommand{\eps}{\varepsilon}
\newcommand{\dd}{\,\mathrm{d}}
\begin{document}
\section*{UnsolvedMath ID: 30003657}
\textbf{Source problem code:} catalogue entry
\subsection{Definitions and mathematical statement}
The source does not define a \emph{core Conceptual Foundations of Mathematics}, specify the philosophical approaches under consideration, or give a criterion for a concept to be acceptable within an approach. With those qualifications unresolved, the supplied question is:
\begin{quote}
Can there be a core Conceptual Foundations of Mathematics consisting of concepts acceptable across several philosophical approaches?
\end{quote}
The cleaned statement asks whether such a common core can be developed across several \emph{major} philosophical approaches. The original statement does not include the restriction to major approaches. Neither formulation specifies whether the approaches are fixed in advance or may be selected together with the core, what substantive content the core must contain, or what would establish feasibility. Consequently, the supplied evidence supports this conceptual feasibility question but does not determine a precise mathematical existence statement.
\subsection{Short English statement}
Can mathematics have a common conceptual foundation acceptable across several philosophical approaches? The intended approaches, the required content of the common core, and the meaning of acceptability remain unspecified.
\subsection{Sources}
\begin{itemize}
\item Local UnsolvedMath registry, permanent ID \texttt{30003657}, title \emph{Shared Conceptual Foundations of Mathematics}, source problem code \texttt{catalogue entry}, collection \emph{Oberwolfach Reports Open Problems}.
\item Supplied source citation: \emph{Mathematical Logic: Proof Theory, Constructive Mathematics} (2018), Oberwolfach Reports, DOI \texttt{10.4171/owr/2017/53}, pp.~3121--3183. \url{https://doi.org/10.4171/owr/2017/53}. The background describes these reports as documenting mathematical workshops and their open-problem discussions.
\item Supplied literature source and attribution context: ``Oberwolfach Report 53, Aczel talk,'' \url{https://doi.org/10.4171/owr/2017/53}. The registry lists proposed year 2017 but supplies no proposer in its proposer field; attribution to a specific proposer is therefore unresolved.
\item Literature assessment dated \texttt{2026-08-21}: the supplied curation labels the problem open, describes the proposal as a feasibility question, and reports finding no subsequent mathematical resolution. This is supplied dated curation, not an independently verified literature assessment. The supplied background requests independent verification of time-sensitive status and cited claims.
\end{itemize}
\end{document}
